\chapter*{Resumen}

El Modelo Estándar (ME) y la teoría electrodébil (EW, por sus siglas en inglés), así como las correcciones radiativas al ángulo de Weinberg ($\sin^2\theta_W$), predicen que los neutrinos, aunque eléctricamente neutros, poseen propiedades electromagnéticas cuantificables. Específicamente, una de estas propiedades es el radio de carga del neutrino (NCR, por sus siglas en inglés). Presentamos cómo diferentes configuraciones experimentales que involucran antineutrinos de reactor interactuando vía Dispersión Elástica Coherente Neutrino-Núcleo (CE$\nu$NS, por sus siglas en inglés) en argón líquido (LAr) pueden obtener resultados diferentes dependiendo de los valores de $\sin^2\theta_W$ y del NCR, cuantificando así la sensibilidad de dichos experimentos a estos parámetros del ME. Los resultados muestran la sensibilidad proyectada de experimentos basados en reactor cerca de la Ciudad de México para determinar el ángulo de Weinberg y el NCR del neutrino electrónico. Encontramos que la precisión proyectada en ambos parámetros escala de forma cercana a $1/\sqrt{N}$ con el tiempo de exposición y mejora en casi un orden de magnitud a lo largo de las masas de detector consideradas, de 10 a 200 kg, mientras que la inclusión de un factor de quenching realista para el argón, en lugar de una respuesta unitaria idealizada, degrada la sensibilidad a ambos parámetros por un factor de dos a tres debido a la pérdida de eventos de retroceso de baja energía por debajo del umbral. Incluso bajo este escenario realista, se proyecta que un detector de 200 kg expuesto durante 200 días restrinja $\sin^2\theta_W$ al nivel de $1$-$2\%$ y el radio de carga del neutrino electrónico a $\langle r_{\nu_e}^2\rangle \sim 2\times10^{-33}\ \text{cm}^2$, una precisión competitiva con las determinaciones de baja energía existentes del ángulo de mezcla débil y un objetivo concreto y comprobable para la estructura electromagnética del neutrino electrónico utilizando una fuente y una tecnología de detección ya disponibles a escala de un laboratorio nacional.
